\bookchapter{Глава 14: Менеджеры паролей — специальная тема}

\chapterintro{В цифровую эпоху пароль — это ключ от вашего дома. Но что делать, если ключей слишком много и они слишком сложные? Менеджер паролей — это ваш «универсальный хранитель ключей».}

\booksection{Зачем вам нужен менеджер паролей?}

У современного человека в среднем десятки, а то и сотни учётных записей в интернете. Если использовать один и тот же пароль для всех аккаунтов, то как только один сайт будет взломан — под угрозой окажутся все. Если использовать разные сложные пароли для каждого аккаунта — их просто невозможно запомнить.

Менеджер паролей решает эту дилемму: вам нужно запомнить только один мастер-пароль, а он запомнит все остальные.

\booksection{Как работает менеджер паролей}

\begin{enumerate}
  \item \textbf{Шифрованное хранение}: все пароли хранятся в зашифрованном виде с использованием стойкого шифрования, и только вы можете их расшифровать.
  \item \textbf{Мастер-пароль}: вы задаёте один мастер-пароль для разблокировки хранилища паролей. Сам мастер-пароль не хранится на сервере.
  \item \textbf{Автозаполнение}: на странице входа на сайт или в приложение менеджер паролей автоматически подставляет логин и пароль.
  \item \textbf{Генерация паролей}: может создавать случайные, сложные пароли для повышения безопасности.
\end{enumerate}

\booksection{Введение в принципы шифрования}

Понимание принципов шифрования, лежащих в основе менеджеров паролей, поможет вам оценить, заслуживает ли сервис доверия. Ниже — объяснение ключевых концепций простым языком.

\booksubsection{Шифрование AES-256}

AES (Advanced Encryption Standard) — симметричный алгоритм шифрования, одобренный Национальным институтом стандартов и технологий США (NIST), широко используется банками, военными и правительственными структурами. «256» означает длину ключа в 256 бит — теоретически для взлома методом перебора потребуется время, превышающее возраст Вселенной. Все основные менеджеры паролей (1Password, Bitwarden, KeePass и др.) используют AES-256 для шифрования вашего хранилища паролей.

\booksubsection{Архитектура нулевого знания (Zero-Knowledge Architecture)}

«Нулевое знание» не означает, что сервис ничего не хранит, — это означает, что\textbf{шифрование и расшифровка происходят локально на вашем устройстве, а сервер хранит только зашифрованный текст}. Конкретный процесс:

\begin{enumerate}
  \item Вы вводите мастер-пароль в клиенте.
  \item Клиент извлекает ключ шифрования из мастер-пароля (это происходит на вашем телефоне/компьютере).
  \item На сервер загружается только зашифрованный текст.
  \item При следующем входе сервер отправляет зашифрованный текст обратно, и клиент снова расшифровывает его с помощью мастер-пароля.
\end{enumerate}

Ключевой момент:\textbf{сам мастер-пароль никогда не покидает ваше устройство}, и на сервере его тоже нет. Поэтому даже если сервис будет взломан, хакер получит лишь стопку нечитаемого шифротекста — если только он не сможет угадать ваш мастер-пароль.

\begin{tcbinfo}{Почему «нулевое знание» означает, что сервис тоже не видит ваши пароли?}
Потому что ключ расшифровки извлекается из вашего мастер-пароля, а мастер-пароль существует только в вашей голове и используется только на вашем локальном устройстве. На серверах провайдера нет ни мастер-пароля, ни ключа расшифровки — только шифротекст. Именно поэтому такие надёжные сервисы, как Bitwarden и 1Password, прямо заявляют: «Мы тоже не можем видеть ваши пароли» — не потому, что не хотят, а потому, что технически не могут.

\end{tcbinfo}

\booksubsection{Функции формирования ключа: PBKDF2 и Argon2}

Мастер-пароль нельзя использовать напрямую как ключ шифрования — он должен пройти через «функцию формирования ключа». Эта функция\textbf{намеренно замедляет вычисления}, делая стоимость brute-force атаки непомерно высокой.

\begin{itemize}
  \item \textbf{PBKDF2}: проверенный временем стандарт, широко используется в iCloud Keychain, Bitwarden, LastPass и других. Замедляет взлом за счёт многократного итерированного хеширования (обычно 100 000+ итераций).
  \item \textbf{Argon2}: победитель конкурса хеш-функций для паролей 2015 года, более современный, чем PBKDF2. Помимо «вычислительной сложности» вводит «потребление памяти» — что затрудняет ускорение взлома даже на специализированном оборудовании, таком как GPU/ASIC. Уже поддерживается в Bitwarden, KeePassXC и других.
\end{itemize}

Чем больше итераций, тем выше безопасность, но и разблокировка будет медленнее. Надёжные сервисы открыто публикуют используемые алгоритмы формирования ключа и количество итераций для аудита сообществом безопасности.

\booksubsection{Предостерегающий пример: утечка данных LastPass в 2022 году}

В августе 2022 года ветеран среди менеджеров паролей LastPass сообщил о взломе своей среды разработки; в ноябре того же года было дополнительно раскрыто, что хакеры получили резервные копии пользовательских хранилищ паролей (в зашифрованном виде). Ключевые уроки этого инцидента:

\begin{itemize}
  \item \textbf{Архитектура нулевого знания устояла}: поскольку само хранилище паролей зашифровано, хакеры получили шифротекст, но не смогли его расшифровать — при условии, что мастер-пароль пользователя был достаточно надёжным.
  \item \textbf{Пользователи со слабыми мастер-паролями — в наибольшей опасности}: хакеры могут офлайн перебирать слабые мастер-пароли, поэтому пользователи со слабыми или повторно используемыми мастер-паролями подвергаются реальному риску.
  \item \textbf{Метаданные тоже утекают}: хотя содержимое паролей зашифровано, метаданные, такие как URL-адреса сайтов и заметки, могут быть не зашифрованы, что раскрывает, какие сервисы использует пользователь.
\end{itemize}

\begin{tcbwarn}{Чему нас учит инцидент с LastPass}
Во-первых, мастер-пароль должен быть длинным и уникальным (рекомендуется не менее 16 символов, включая заглавные и строчные буквы, цифры и символы). Во-вторых, отдавайте предпочтение сервисам, поддерживающим Argon2id с высоким числом итераций. В-третьих, для наиболее ценных учётных записей (банк, почта) обязательно включите отдельную двухфакторную аутентификацию, не полагайтесь только на менеджер паролей.

\end{tcbwarn}

\booksection{Рекомендации популярных менеджеров паролей}

\booksubsection{Встроенные в систему:}

\begin{itemize}
  \item \textbf{iCloud Keychain (Связка ключей iCloud)}: лучший выбор для пользователей экосистемы Apple. Бесшовная интеграция с Safari и приложениями, бесплатно, безопасно. Недостаток — неудобство кросс-платформенного использования.
  \item \textbf{Google Password Manager}: удобно для пользователей Android. Интеграция с Chrome и системой Android, бесплатно.
  \item \textbf{Huawei Password Vault}: встроено в HarmonyOS, подходит для пользователей Huawei.
\end{itemize}

\booksubsection{Сторонние профессиональные инструменты:}

\begin{itemize}
  \item \textbf{Bitwarden}: с открытым исходным кодом, бесплатный, мощный, поддерживает все платформы. Настоятельно рекомендуется пользователям, ценящим приватность.
  \item \textbf{1Password}: красивый интерфейс, богатый функционал, подходит для семей и команд. Платный.
  \item \textbf{LastPass}: ветеран среди менеджеров паролей, функционал бесплатной версии ограничен.
  \item \textbf{KeePass}: полностью локальное хранение, с открытым исходным кодом, бесплатный. Подходит для гиков.
\end{itemize}

\booksection{Реальная оценка надёжности паролей}

Менеджеры паролей обычно показывают индикатор «надёжности пароля» (красный/жёлтый/зелёный), но эта оценка часто учитывает только длину и типы символов и не является всесторонней. Чтобы действительно определить, безопасен ли пароль, нужно учитывать «энтропию» и то, был ли он уже скомпрометирован.

\booksubsection{Энтропия: научная мера надёжности пароля}

«Энтропия» измеряет неопределённость пароля, единица измерения — биты. Чем выше энтропия, тем больше попыток потребуется для brute-force атаки.

\begin{itemize}
  \item \textbf{«123456»}: крайне низкая энтропия, взламывается мгновенно.
  \item \textbf{«P@ssw0rd»}: выглядит сложным, но присутствует во всех списках словарных атак, энтропия всё ещё низкая.
  \item \textbf{«correct horse battery staple»}: четыре случайных слова, энтропия около 44 бит, сложность взлома намного выше, чем у «P@ssw0rd».
  \item \textbf{16 случайных символов (заглавные + строчные + цифры + символы)}: энтропия около 100 бит, при текущих вычислительных мощностях brute-force атака невозможна.
\end{itemize}

Вывод:\textbf{«длина + случайность» важнее «сложных символов»}. Именно поэтому случайные длинные пароли, генерируемые менеджером паролей, действительно безопасны.

\booksubsection{HaveIBeenPwned: проверьте, не утёк ли ваш пароль}

Исследователь безопасности Трой Хант поддерживает сайт\href{https://haveibeenpwned.com}{haveibeenpwned.com}, который собирает учётные записи и пароли, утёкшие в ходе крупных утечек данных. Вы можете:

\begin{enumerate}
  \item Ввести email, чтобы проверить, был ли ваш аккаунт раскрыт в какой-либо утечке.
  \item Использовать функцию Pwned Passwords для проверки, есть ли конкретный пароль в известных базах утёкших паролей (менеджеры паролей и браузеры уже встроили эту проверку).
\end{enumerate}

Если пароль найден в базе утёкших,\textbf{каким бы сложным он ни выглядел, его необходимо немедленно заменить}— потому что словарные атаки в первую очередь пробуют эти известные пароли.

\begin{tcbinfo}{Принцип запроса k-Anonymity}
Беспокоитесь, что ввод пароля на HaveIBeenPwned может снова привести к утечке? Не стоит. Современные запросы используют протокол «k-анонимности»: браузер отправляет на сервер только первые 5 символов хеша пароля, сервер возвращает список всех хешей, начинающихся с этих 5 символов, а сравнение происходит локально. То есть сервер никогда не видит полный хеш вашего пароля. Именно так работают функции «обнаружения утёкших паролей» в 1Password и Bitwarden.

\end{tcbinfo}

\booksubsection{Встроенные функции мониторинга в менеджерах паролей}

\begin{itemize}
  \item \textbf{1Password Watchtower}: в реальном времени отслеживает слабые пароли, повторяющиеся пароли, утёкшие пароли, просроченные пароли, сайты без 2FA и отображает их с цветовой кодировкой (красный/оранжевый/жёлтый).
  \item \textbf{Bitwarden Vault Health}: бесплатная версия показывает слабые и повторяющиеся пароли; платная версия обнаруживает утёкшие пароли (на основе HaveIBeenPwned).
  \item \textbf{«Рекомендации по безопасности» Связки ключей iCloud}: Настройки iOS Пароли Рекомендации по безопасности — покажет, какие пароли найдены в известных утечках и какие используются повторно.
  \item \textbf{«Проверка паролей» Google Password Manager}: встроена в Chrome, позволяет одним кликом проверить утёкшие, слабые и повторяющиеся пароли.
\end{itemize}

\begin{tcbpractice}{Проверьте прямо сейчас: безопасны ли ваши пароли}
\begin{enumerate}
  \item Откройте haveibeenpwned.com, введите ваш основной email и посмотрите, в каких утечках он фигурирует.
  \item Пользователи iPhone: Настройки Пароли Рекомендации по безопасности, обработайте по очереди все красные предупреждения.
  \item Пользователи Chrome: Настройки Пароли Проверить пароли, измените согласно рекомендациям.
  \item Пользователи 1Password/Bitwarden: откройте Watchtower/Health Reports, в первую очередь обработайте пункты «Утёкшие» и «Повторяющиеся».
\end{enumerate}

\end{tcbpractice}

\booksection{Passkey: будущее беспарольного входа}

В 2024–2026 годах технологии паролей переживают революцию —\textbf{Passkey (ключ доступа)}обещает полностью вытеснить традиционные пароли. Apple, Google и Microsoft совместно продвигают эту технологию при поддержке альянса FIDO и Консорциума Всемирной паутины (W3C) — это крупнейшая реформа способов входа за последние десятилетия.

\booksubsection{Что такое Passkey?}

Passkey основан на стандартах FIDO2 / WebAuthn и заменяет пароли\textbf{парой ключей на устройстве}. Принцип можно упрощённо представить так:

\begin{itemize}
  \item Ваш телефон генерирует пару ключей:\textbf{открытый ключ}(регистрируется на сайте) +\textbf{закрытый ключ}(навсегда остаётся в защищённом чипе вашего устройства).
  \item При входе сайт отправляет «вызов», ваш телефон подписывает его закрытым ключом, а сайт проверяет подпись открытым ключом.
  \item Перед подписанием требуется подтверждение личности через биометрию (Face ID / отпечаток пальца) или PIN-код устройства.
  \item В течение всего процесса\textbf{ни один пароль не передаётся по сети}.
\end{itemize}

\booksubsection{Passkey vs традиционные пароли}

\begin{center}
\begin{tabularx}{\linewidth}{|X|X|X|}
\hline
\textbf{Аспект} & \textbf{Традиционные пароли} & \textbf{Passkey} \tabularnewline
\hline
Защита от фишинга & Слабая, фишинговые сайты могут украсть пароль & Врождённая защита от фишинга, закрытый ключ подписывает только для подлинного домена \tabularnewline
Влияние утечки на сервере & Хеш пароля может быть взломан & Сервер хранит только открытый ключ, утечка не наносит вреда \tabularnewline
Удобство & Нужно запоминать/вводить/управлять & Вход за секунду по Face ID/отпечатку пальца \tabularnewline
Атаки повторного воспроизведения & Могут быть перехвачены и воспроизведены & Каждая подпись одноразовая, невозможно воспроизвести \tabularnewline
\hline
\end{tabularx}
\end{center}

\booksubsection{Поддержка на разных платформах}

\begin{itemize}
  \item \textbf{Apple iCloud Keychain Passkey}: нативная поддержка в iOS 16+ / macOS Ventura+, Passkey синхронизируются между всеми устройствами с одним Apple ID, при входе на сайт с нового устройства не нужно создавать заново.
  \item \textbf{Google Password Manager Passkey}: нативная поддержка в Android + Chrome, Passkey хранятся в аккаунте Google и синхронизируются между устройствами Android; с 2023 года личные аккаунты Google поддерживают вход по Passkey.
  \item \textbf{Microsoft}: Windows 11 / Edge уже поддерживают Passkey; сервисы Microsoft (учётная запись, Outlook, Office и др.) постепенно открывают вход по Passkey, управление доступно в Microsoft Authenticator.
  \item \textbf{1Password / Bitwarden}: уже поддерживают создание, хранение и кросс-платформенную синхронизацию Passkey, позволяя не ограничиваться одной экосистемой.
\end{itemize}

\begin{tcbinfo}{Кросс-платформенная синхронизация — ключевой прорыв}
Главной проблемой ранних Passkey была «привязка к одному устройству» — при смене телефона приходилось создавать всё заново. С 2023 года Apple, Google и Microsoft реализовали облачную синхронизацию Passkey: Связка ключей iCloud синхронизирует между устройствами Apple, Google Password Manager — между устройствами Android. Сторонние решения, такие как 1Password, обеспечивают кросс-экосистемную синхронизацию. Это означает, что Passkey уже пригодны для практической замены паролей.

\end{tcbinfo}

\booksubsection{Как создать и использовать Passkey на телефоне}

\textbf{Для iOS (iPhone):}

\begin{enumerate}
  \item На странице безопасности аккаунта сайта, поддерживающего Passkey (например, Google, Microsoft, Amazon, eBay, PayPal), выберите «Добавить ключ доступа / Add passkey».
  \item Система покажет окно проверки Face ID / Touch ID, после подтверждения Passkey автоматически создаётся и сохраняется в Связку ключей iCloud.
  \item При следующем входе на этот сайт введите имя пользователя, выберите «Войти с помощью iPhone», проведите Face ID — и готово.
  \item Управление: Настройки Пароли найдите нужный сайт, там будет запись Passkey.
\end{enumerate}

\textbf{Для Android:}

\begin{enumerate}
  \item На странице безопасности аккаунта сайта, поддерживающего Passkey, выберите «Добавить ключ доступа».
  \item Всплывёт окно Google Password Manager выберите ваш аккаунт Google подтвердите блокировкой экрана (отпечаток/графический ключ/PIN).
  \item Passkey автоматически сохраняется в Google Password Manager и синхронизируется между устройствами Android с тем же аккаунтом.
  \item При следующем входе выберите «Использовать Passkey», разблокируйте экран — и готово.
  \item Управление: Настройки Android Пароли, ключи и автозаполнение Менеджер паролей google.com.
\end{enumerate}

\booksubsection{Кросс-платформенный вход с Passkey}

Если вы входите на сайте на компьютере с Windows, а Passkey есть только на iPhone, можно авторизоваться через QR-код:

\begin{enumerate}
  \item На странице входа выберите «Использовать Passkey».
  \item Выберите «Использовать другое устройство».
  \item Компьютер покажет QR-код, отсканируйте его iPhone.
  \item iPhone запросит Face ID, после подтверждения вход на сайте будет выполнен на компьютере.
\end{enumerate}

Этот процесс использует двойное шифрование через Bluetooth + интернет, закрытый ключ Passkey не покидает ваш iPhone.

\begin{tcbpractice}{Попробуйте прямо сейчас: испытайте беспарольный вход}
\begin{enumerate}
  \item Войдите в аккаунт Google Безопасность Ключи доступа Добавить ключ доступа, испытайте вход по Face ID/отпечатку пальца.
  \item Войдите в аккаунт Microsoft Безопасность Дополнительные параметры безопасности Добавить ключ доступа.
  \item Попробуйте Passkey на сервисах, уже поддерживающих эту функцию: Taobao, JD.com, eBay, PayPal и др.
  \item Сохраните традиционный пароль как резервный и постепенно переводите наиболее ценные учётные записи на Passkey.
\end{enumerate}

\end{tcbpractice}

\begin{tcbwarn}{Рекомендации на переходный период}
Passkey ещё не полностью заменил пароли — многие сайты находятся в стадии «Passkey + пароль». Рекомендация: для наиболее ценных учётных записей включайте Passkey в первую очередь, но\textbf{не удаляйте сразу традиционный пароль}— оставьте его как резервный; продолжайте использовать менеджер паролей, так как не все сайты поддерживают Passkey. Рассматривайте Passkey как «мощное дополнение к менеджеру паролей», а не как «немедленную замену».

\end{tcbwarn}

\booksection{Аппаратные ключи безопасности}

Если ваши требования к безопасности выше, чем у обычного пользователя — например, вы владеете большим количеством криптовалюты, являетесь руководителем/журналистом/активистом или защищаете ключевые корпоративные учётные записи — аппаратный ключ безопасности является более надёжным выбором, чем SMS-коды подтверждения.

\booksubsection{Что такое аппаратный ключ безопасности?}

Аппаратный ключ — это физическое устройство (похожее на USB-флешку), внутри которого защищённый чип хранит ваш закрытый ключ. Для проверки требуется физический контакт или NFC. Представительные продукты:

\begin{itemize}
  \item \textbf{YubiKey}(компания Yubico): самый популярный, поддерживает FIDO2/WebAuthn, U2F, OTP, PIV, OpenPGP, интерфейсы USB-C / USB-A / NFC и другие. Серия 5 стоит 50–70 долларов.
  \item \textbf{Google Titan}: собственный продукт Google, USB-C + NFC / Bluetooth, около 30–40 долларов.
  \item \textbf{OnlyKey}: альтернатива с открытым исходным кодом, поддерживает несколько протоколов.
  \item \textbf{Китайские аналоги}: например, Feitian ePasskey, более доступная цена.
\end{itemize}

\booksubsection{Почему аппаратные ключи безопаснее SMS-кодов?}

\begin{itemize}
  \item \textbf{Защита от фишинга}: как и Passkey, аппаратный ключ отвечает только на подлинный домен, фишинговые сайты не могут выманить код подтверждения.
  \item \textbf{Защита от SIM-подмены}: SMS-коды могут быть перехвачены атакой SIM-swap, аппаратному ключу не нужен номер телефона.
  \item \textbf{Защита от атак «человек посередине»}: протокол FIDO2 включает значение вызова в каждую подпись, что делает невозможным повторное воспроизведение.
  \item \textbf{Физическая изоляция}: закрытый ключ навсегда остаётся внутри аппаратного чипа, вредоносное ПО не может его прочитать, даже если ваш компьютер взломан.
\end{itemize}

\booksubsection{Сценарии применения}

\begin{itemize}
  \item Аккаунты на криптовалютных биржах, восстановление кошельков.
  \item Корпоративный SSO, удалённый доступ VPN, SSH-вход на серверы.
  \item «Основные учётные записи» — почта, GitHub, социальные сети — взлом которых позволяет сбросить все остальные аккаунты.
  \item Лица с высокой чувствительностью: журналисты, юристы, руководители.
\end{itemize}

\booksubsection{Как настроить аппаратный ключ как второй фактор}

На примере аккаунта Google + YubiKey:

\begin{enumerate}
  \item Купите YubiKey (рекомендуется сразу две штуки: одну для повседневного использования, другую как резервную, хранить в разных местах).
  \item Войдите на myaccount.google.com Безопасность Двухэтапная аутентификация.
  \item Выберите «Добавить ключ безопасности» вставьте YubiKey в USB-C порт или поднесите к NFC-зоне телефона.
  \item Следуя подсказкам, коснитесь металлического контакта YubiKey для завершения регистрации.
  \item Рекомендуется зарегистрировать второй резервный ключ.
  \item В настройках «Двухэтапной аутентификации» установите аппаратный ключ как предпочтительный способ проверки, SMS-коды — как резервный.
\end{enumerate}

Основные поддерживаемые сайты: Google, GitHub, Microsoft, Apple ID, Facebook, Twitter, Dropbox, GitLab, AWS, Coinbase, Binance и др.

\begin{tcbinfo}{Стратегия двух ключей}
Самый большой риск аппаратного ключа — «потеря». Рекомендация: (1) зарегистрируйте два разных ключа одновременно; (2) один носите с собой (на брелоке), другой храните в домашнем сейфе или банковской ячейке; (3) некоторые сайты позволяют распечатать «резервные коды восстановления» и хранить их офлайн. Так вы не будете навсегда заблокированы вне аккаунта, даже если потеряете основной ключ.

\end{tcbinfo}

\begin{tcbwarn}{Будьте осторожны с каналом покупки}
Обязательно покупайте в официальном магазине или у авторизованного дилера, не покупайте подержанные ключи или ключи неизвестного происхождения — теоретически предварительно настроенное оборудование может содержать бэкдор. Прошивку YubiKey нельзя обновить после выпуска, но Yubico постоянно публикует бюллетени безопасности — следите за их официальным сайтом.

\end{tcbwarn}

\booksection{Совместное использование паролей в семье и команде}

В семейном сценарии пароли от Netflix, домашнего Wi-Fi, роутера и умного замка должны использоваться всей семьёй; в командном сценарии серверы, базы данных и аккаунты в соцсетях требуют совместной работы нескольких человек.\textbf{Отправлять пароли через WeChat или записывать их в общих документах — крайне опасная практика}: как только история чата или права доступа к документу будут скомпрометированы, пароль окажется раскрыт.

\booksubsection{«Общие хранилища» (Vault) в 1Password}

1Password — эталонный продукт для семейного и командного обмена. Его ключевой механизм —\textbf{несколько хранилищ}:

\begin{itemize}
  \item \textbf{Личное хранилище}: только вы можете видеть, например, банковские и почтовые учётные записи.
  \item \textbf{Общее хранилище}: могут видеть члены семьи или команды, например, Netflix, домашний Wi-Fi, общие кредитные карты.
  \item \textbf{Разграничение прав}: в командной версии можно настроить для некоторых участников «только просмотр» или «только просмотр без редактирования».
\end{itemize}

Членам семьи или команды не нужно делиться мастер-паролем — каждый входит под своей учётной записью, но может видеть содержимое авторизованных общих хранилищ. При изменении пароля все участники получают обновление в реальном времени.

\booksubsection{Функция «Организации» в Bitwarden}

Бесплатная версия Bitwarden поддерживает личный обмен паролями (через функцию «Send» — ограниченное по времени предоставление одной записи); платная версия позволяет создавать «Организации» для совместного доступа к определённым записям с членами семьи или коллегами с детальным разграничением прав. Bitwarden также предлагает вариант самостоятельного хостинга, подходящий для предприятий с высокими требованиями к суверенитету данных.

\booksubsection{Семейный сценарий: безопасный способ делиться паролями Netflix, Wi-Fi и т.д.}

\begin{itemize}
  \item \textbf{Netflix / видеосервисы}: храните пароль для входа в «Общем хранилище» семейной версии 1Password, каждый член семьи входит самостоятельно — не нужно отправлять в чат.
  \item \textbf{Пароль домашнего Wi-Fi}: между устройствами iOS можно напрямую поделиться паролем Wi-Fi через AirDrop с гостем, не произнося его вслух; Android 12+ также имеет функцию «Поделиться Wi-Fi поблизости».
  \item \textbf{Умные замки и умная бытовая техника}: по возможности используйте функцию «Члены семьи» в приложении производителя для создания отдельного аккаунта для каждого человека, а не используйте общий административный аккаунт.
  \item \textbf{Панель управления домашним роутером}: пароль администратора храните в общем хранилище 1Password, для гостей создайте отдельный SSID гостевого Wi-Fi.
\end{itemize}

\begin{tcbinfo}{Соотношение цены и качества семейной версии 1Password}
1Password Families стоит около 4,99 доллара в месяц и поддерживает до 5 членов семьи, каждый имеет независимое личное хранилище + 4 общих хранилища. В пересчёте на человека это менее 1 доллара в месяц — лучший вариант для семейного обмена. Семейная версия Bitwarden стоит 10 долларов в год на 6 человек — ещё выгоднее, но интерфейс немного уступает.

\end{tcbinfo}

\booksubsection{Командный сценарий: обмен паролями серверов и баз данных}

\begin{itemize}
  \item \textbf{1Password Teams / Business}: можно создавать несколько хранилищ по отделам и проектным группам, детальный контроль прав доступа, все действия подлежат аудиту.
  \item \textbf{Bitwarden Teams / Enterprise}: решение с открытым исходным кодом, возможность самостоятельного хостинга, корпоративная версия поддерживает интеграцию SSO и синхронизацию каталогов.
  \item \textbf{Принцип разделения обязанностей}: пароль администратора не должен быть известен только одному человеку — назначьте 2–3 «экстренных контакта», чтобы обеспечить непрерывность работы при увольнении или отпуске ключевого сотрудника.
  \item \textbf{Процедура при увольнении}: при увольнении сотрудника немедленно отзовите его доступ к менеджеру паролей и при необходимости смените все высокочувствительные пароли, к которым у него был доступ.
\end{itemize}

\booksubsection{Экстренный доступ (Emergency Access): на случай непредвиденных обстоятельств}

Представьте сценарий: вы единолично управляете всеми важными аккаунтами семьи, но попадаете в больницу или, к несчастью, умираете — и семья не может войти в банковские, почтовые и страховые учётные записи.\textbf{Функция «Экстренный доступ»}создана именно для этого — она позволяет доверенному члену семьи или другу получить доступ к вашему хранилищу паролей, если вы не можете ответить.

\textbf{Настройка экстренного доступа в Bitwarden:}

\begin{enumerate}
  \item Войдите в Bitwarden Web Vault Настройки Экстренный доступ.
  \item Добавьте email доверенного контакта и пригласите его принять.
  \item Установите «период ожидания» (рекомендуется 7–14 дней): временной интервал от запроса доступа до вашего подтверждения.
  \item Если вы не ответите в течение периода ожидания (отказ или бездействие), по истечении срока контакт автоматически получит доступ.
  \item Вы в любое время можете вручную отклонить запрос или отозвать право экстренного доступа.
\end{enumerate}

\textbf{Экстренный доступ в 1Password:}

\begin{enumerate}
  \item 1Password не называет это «экстренным доступом» напрямую, но это реализуется через функцию «Семейная/командная организация» — добавьте семью в семейную версию и поместите часть записей в «Общее хранилище».
  \item Для пользователей персональной версии рекомендуется сообщить мастер-пароль высокодоверенному члену семьи или записать его в запечатанный конверт и хранить в сейфе.
  \item 1Password также предоставляет «Код восстановления» — обязательно распечатайте и сохраните его при включении двухфакторной аутентификации.
\end{enumerate}

\begin{tcbpractice}{Настройте прямо сейчас: оставьте семье «запасной ключ»}
\begin{enumerate}
  \item Выберите доверенного члена семьи в качестве «экстренного контакта».
  \item Настройте экстренный доступ в менеджере паролей (Bitwarden) или добавьте в семейное общее хранилище (1Password).
  \item Запишите мастер-пароль + код восстановления на бумаге, запечатайте в конверт и храните в домашнем сейфе или передайте доверенному члену семьи.
  \item Раз в полгода пересматривайте, остаётся ли экстренный контакт подходящим.
\end{enumerate}

\end{tcbpractice}

\booksection{Лучшие практики использования менеджера паролей}

\begin{itemize}
  \item \textbf{Установите надёжный мастер-пароль}: мастер-пароль — это последняя линия обороны; обязательно задайте длинный и сложный пароль и запомните его.
  \item \textbf{Включите двухфакторную аутентификацию}: включите 2FA для аккаунта менеджера паролей, чтобы добавить дополнительный уровень защиты.
  \item \textbf{Регулярно проверяйте слабые пароли}: большинство менеджеров паролей имеют функцию «аудита безопасности», которая помогает найти слабые и повторяющиеся пароли.
  \item \textbf{Не забывайте мастер-пароль}: мастер-пароль невозможно восстановить; забыть его — значит потерять все пароли. Рекомендуется записать его на бумаге и хранить в безопасном месте.
  \item \textbf{Экстренный доступ}: настройте экстренного контакта, чтобы в случае вашей недоступности доверенное лицо могло получить доступ к вашим паролям.
\end{itemize}

\begin{tcbwarn}{Важное напоминание}
Менеджер паролей удобен, но несёт и риски: если мастер-пароль будет раскрыт, все пароли окажутся под угрозой. Поэтому мастер-пароль ни в коем случае не должен совпадать с другими паролями и не должен быть записан в легкодоступном месте. Кроме того, рекомендуется выбирать сервис с хорошей репутацией в области безопасности.

\end{tcbwarn}

\begin{tcbpractice}{Попробуйте прямо сейчас: начните использовать менеджер паролей}
\begin{enumerate}
  \item Выберите менеджер паролей (рекомендуется Bitwarden или встроенный в систему).
  \item Зарегистрируйте аккаунт и задайте надёжный мастер-пароль.
  \item Импортируйте существующие пароли (если есть) или начните добавлять новые.
  \item Установите расширение для браузера или мобильное приложение, чтобы испытать функцию автозаполнения.
  \item Проверьте отчёт аудита безопасности и найдите пароли, которые нужно улучшить.
\end{enumerate}

\end{tcbpractice}
